\documentclass[11pt]{article}
\usepackage[margin=0.95in]{geometry}
\usepackage[T1]{fontenc}
\usepackage{lmodern}
\usepackage[hidelinks]{hyperref}
\usepackage{booktabs}
\usepackage{microtype}
\newcommand{\code}[1]{\texttt{\detokenize{#1}}}

\title{SLURM-Only Memory Enforcement for Tactix Experiments}
\author{Tactix Development Project}
\date{27 August 2026}

\begin{document}
\emergencystretch=2em
\maketitle

\section*{Outcome}
The maintained benchmark infrastructure no longer uses \code{RLIMIT_AS}.
SLURM's task cgroup is now the sole memory-enforcement authority.  Allocation
is expressed as 8~GiB per physical core; on the current SMT-2 nodes a strict
16-GiB task therefore requests two physical/four logical CPUs, while a targeted
timing task requests six physical/twelve logical CPUs and exactly 48,000~MiB.
No solver semantics changed.

Per project-owner direction, the current rank-1--40 campaign will be preserved
as evidence under its original dual-limit contract.  A fresh 80-arm screen will
then rerun the same accepted Tactix and current-Kissat pins with no address-space
limit and will become the selection authority.

\section{Why the old jobs failed below 16 GiB RSS}
Campaign \code{20260826-2340-tierd-ranks1-40-current-screen} requested
\code{--mem=48000M} from SLURM, but its copied runner separately imposed
\code{RLIMIT_AS=16 GiB}.  Tactix ranks 5, 13, and 14 reported
\code{std::bad_alloc} at peak RSS values of 13,970,768, 14,189,420, and
13,342,516~KiB.  These were not 48,000-MiB cgroup OOM events.

An address-space limit counts virtual mappings.  A capacity-growing vector or
hash table may allocate replacement storage while its old storage remains live;
that request can exceed the virtual-space ceiling even though resident high-water
is much lower.  The old screen therefore cannot support a memory-competitiveness
decision.

\section{Cluster enforcement evidence}
An authoritative \code{scontrol show config} query at
\code{2026-08-27T00:50:26-04:00} reported:
\begin{itemize}
\item \code{ProctrackType=proctrack/cgroup} and
  \code{TaskPlugin=task/cgroup};
\item \code{SelectTypeParameters=CR_CPU_MEMORY}; and
\item \code{ConstrainRAMSpace=yes}, \code{AllowedRAMSpace=100.0\%}, and
  \code{AllowedSwapSpace=0.0\%}.
\end{itemize}
Consequently, \code{--mem} is a hard resident-memory task-cgroup ceiling and
swap cannot hide an overage.

\section{Infrastructure and policy changes}
\begin{itemize}
\item \code{tools/bench/run.py} retains its CPU-time backstop but installs no
  memory rlimit.  New SLURM manifests record
  \code{memory_enforcement=slurm_task_cgroup}; local diagnostics explicitly
  record that they have no memory enforcement.
\item The frozen-baseline, Tactix Tier-D, LymphoSAT-attribution, and Phase-4
  diagnostic generators now request total SLURM memory and two logical CPUs per
  physical core.  Their maintained execution paths install no address-space cap.
\item \code{DEVELOPMENT_PLAN.md} now states the physical-core rule and makes the
  strict 16-GiB shape explicit.  It also separates a 48,000-MiB timing run from
  a later strict 16-GiB confirmation.
\item Frozen campaigns and their copied runners remain immutable.  Legacy
  manifest fields are read only for historical interpretation.
\end{itemize}

\section{Validation}
All 76 Release CTest tests pass.  The new memory-policy regression parses each
maintained runner and rejects any \code{setrlimit} call carrying
\code{RLIMIT_AS}; it also guards the SMT-aware allocation formula and plan
language.  Campaign cleanup, Tier-D pinning, and Phase-4 allocation tests pass,
as do Python compilation and all applicable documentation authority checks.
The documentation size checker retains two pre-existing warning-level debts and
reports no error.

No fuzz campaign is required for this increment because neither solver code nor
proof/model semantics changed.

\section{Next action}
After this isolated infrastructure increment is committed and pushed, freeze a
new rank-1--40 screen with the same accepted solvers and immutable inputs.  Run
80 independent proof-producing solver tasks at \code{-c 12}, exactly
\code{--mem=48000M}, 1,200 seconds, no throttle, no node constraint, and no
\code{RLIMIT_AS}; independently validate every model and proof and require zero
certificate residue.  Use its complete aggregate to choose the next gap.  A
genuine Tactix memory solve loss has priority; otherwise proceed to Rank 4 and
then Rank 32.  Confirm any decision-changing 16-GiB claim separately at
\code{-c 4}/\code{--mem=16G}.

\end{document}
